\section{篇章分析}
\label{sec:nlp-analysis-discourse}

把每句话分析正确，仍可能不知道后一句的“他”是谁，也不知道它是在解释、反驳还是补充前一句。\term{篇章分析}（discourse analysis）处理跨句子的联系。图~\ref{fig:nlp-analysis-discourse}用补充段落“小王借到一本书。他读完后，归还了这本书。虽然封面已经磨损，但内容仍然清晰”展示三种不同的联系。“小王”与“他”指同一个人，“一本书”与“这本书”指同一本书，这是身份上的同一；“封面”与书是部分与整体，这是\term{桥接}（bridging），两者相关却不是同一对象；“虽然……但……”连接的是两个陈述，表达让步，联系的端点是内容单元而不是实体。

\begin{figure*}[htbp]
  \centering
  \begin{tikzpicture}[
  x=1mm,y=1mm,
  every node/.append style={font=\figurefont\footnotesize,align=center},
  mention/.style={draw=FigureInput,fill=FigureInput!10,
    rounded corners=1mm,inner sep=1.5mm,minimum height=8mm},
  identity/.style={draw=FigureInput,line width=1.1pt},
  bridge/.style={->,draw=FigureOutput,dashed,line width=1.1pt},
  discourse/.style={->,draw=FigureModel,dash dot,line width=1.1pt}
]
  \node[mention] (wang) at (20,0) {小王};
  \node at (42,0) {借到};
  \node[mention] (book) at (71,0) {一本书};
  \node at (86,0) {。};
  \node[mention] (he) at (20,-20) {他};
  \node at (42,-20) {读完后，};
  \node[circle,draw=FigureInput,line width=1.1pt,
    minimum size=3mm,inner sep=0] (zero) at (67,-20) {};
  \node at (90,-20) {归还了};
  \node[mention] (book2) at (128,-20) {这本书};
  \node at (145,-20) {。};

  \draw[identity] (wang.south) -- (he.north);
  \draw[identity] (wang.south) -- (20,-10) -- (67,-10) -- (zero.north);
  \draw[identity] (book.north) -- (71,9) -- (128,9) -- (book2.north);
  \node[anchor=south] at (113,1) {同一本书};
  \node[anchor=north] at (67,-23) {省略位置};

  \draw[draw=FigureModel,fill=FigureModel!5,rounded corners=1mm]
    (2,-59) rectangle (84,-43);
  \draw[draw=FigureModel,fill=FigureModel!5,rounded corners=1mm]
    (94,-59) rectangle (162,-43);
  \node at (12,-51) {虽然};
  \node[mention,draw=FigureOutput,fill=FigureOutput!10] (cover) at (32,-51) {封面};
  \node at (62,-51) {已经磨损，};
  \node at (128,-51) {但内容仍然清晰。};
  \draw[bridge] (cover.north) -- (32,-35) -- (128,-35) -- (book2.south);
  \node[fill=white,inner sep=1mm] at (102,-35) {PART-OF};
  \draw[discourse] (43,-59) -- (43,-69) -- (128,-69) -- (128,-59);
  \node[fill=white,inner sep=1mm] at (85,-69) {让步：两个陈述};

  \draw[identity] (3,-82) -- (15,-82);
  \node[anchor=west] at (18,-82) {同一身份};
  \draw[bridge] (55,-82) -- (67,-82);
  \node[anchor=west] at (70,-82) {桥接};
  \draw[discourse] (109,-82) -- (121,-82);
  \node[anchor=west] at (124,-82) {篇章关系};
\end{tikzpicture}

  \caption{同一段落中的身份共指、桥接与篇章关系。实线连接同一对象的提及，虚线连接“封面”与书的部分--整体关系，点划线连接篇章单元。空心位置标记不是原文字符，而是“归还”前检测出的省略主语位置，先行项为小王。共指可以形成等价簇，桥接和篇章关系不能并入该簇。}
  \label{fig:nlp-analysis-discourse}
\end{figure*}
\FloatBarrier

三种边的端点和性质不同，输出结构也就不同：身份共指满足传递性，可以合并为等价簇；篇章关系有方向和类型，还可能组织成覆盖全文的树；论证关系同样有方向，但关心的是一个单元是否为另一个单元提供理由。本节据此分为四个小节：共指处理提及的身份，篇章关系处理内容单元之间的语义联系，论证结构处理主张与理由，连贯性则不再输出具体的边，而是给整段文字打分或排定句序。长文档带来的截断和候选筛选问题会影响这四类任务的计算方式，但不构成新的关系类型。

\subsection{提及、先行项与共指}
\label{sec:nlp-analysis-coreference}

\term{提及}（mention）是文本中指向某个对象的一次表达，可以是“机器学习读书会”这样的名词短语，也可以是代词。\term{共指}（coreference）关系连接指向同一对象的提及，所有互相共指的提及构成一个共指簇。\term{先行项}（antecedent）是用来解释当前表达的较早提及。簇与先行项是两种数据结构：簇描述身份，先行项指针描述一条恢复路径；只要把指针看作无向边并取连通分量，指针就能还原为簇，而同一个簇可以对应许多组不同的指针。

有三类现象需要先划出边界。“封面”与书的桥接不合并身份；“活动改期了。这让小王调整了行程”中的“这”指向改期这一事件，属于话语指示，可能没有对应的名词跨度；中文还常省略主语，图中“归还”前的空位需要先检测出位置，再为它选先行项。本节主要讨论名词性提及的身份共指；零代词可视为在字符边界上增加的候选提及，检测与选择先行项分开计分。

共指消解面对的挑战可以归为五项。其一是提及检测：哪些跨度是提及，单例提及（只出现一次的实体）是否要输出。其二是长距离先行项：代词的先行项可能在几句之前，候选数量随文档长度增长。其三是一致性与语义兼容：数、性、生命度要相容，“那家公司”与“它”可以同指，与“他”通常不能。其四是簇级一致性：两两判断各自合理，合起来却可能违背传递性。其五是评价：误合并两个实体与漏连一条边，对不同指标的影响差别很大。

方法的演进大致对应这些挑战。Hobbs算法在句法树上按固定顺序搜索代词的先行项，依靠句法结构和一致性规则处理第三项；提及对模型用分类器判断两两是否同指，可以融合大量特征，但各次判断互相独立；提及排序模型对一个提及的全部候选先行项归一化，只选最好的一个，缓解了独立二分类的冲突；端到端跨度模型再把提及检测并入同一个目标，消除检测与连接之间的流水线误差，并用剪枝和粗到细推断控制候选规模。第~\ref{sec:nlp-analysis-coreference-metrics}小节给出共指评价指标，第~\ref{chap:nlp-information-extraction}章的事件共指沿用这些指标。

\subsubsection{Hobbs句法先行项搜索}
\label{sec:nlp-analysis-hobbs}

在没有标注语料可学的年代，代词消解的主要依据是句法结构：英语中先行项更常出现在主语位置、更常出现在离代词较近的句子里，而且同一小句里的宾语代词一般不与本句主语同指（“John likes him”中him不指John）。Hobbs的朴素算法把这些倾向写成一条在句法树上的固定搜索顺序，按顺序提出候选，用数、性等一致性条件逐个筛选，第一个通过的候选就是结果\scite{nlp-analysis-hobbs1978}。算法的输入是当前句与前文各句的表层句法树，步骤如下：
\begin{enumerate}
  \item 从直接支配代词的NP节点出发。
  \item 向上走到遇到的第一个NP或S节点，记为X，并记下走过的路径p。
  \item 在X之下、路径p左侧的分支中，按从左到右的广度优先顺序遍历，把遇到的、与X之间还隔着某个NP或S节点的NP提为候选。
  \item 若X是本句最高的S节点，转去按由近及远的顺序遍历前文各句的句法树，每棵树内部从左到右广度优先，遇到的NP依次提为候选；否则继续第5步。
  \item 从X继续向上，走到下一个NP或S节点，把它记为新的X，路径记为p。
  \item 若X是NP，且路径p没有经过X直接支配的$\bar N$节点，把X本身提为候选。
  \item 在X之下、路径p左侧的分支中从左到右广度优先遍历，把遇到的NP提为候选。
  \item 若X是S，再遍历X之下路径p右侧的分支，但遇到NP或S节点时不再深入其内部，把遇到的NP提为候选。
  \item 回到第4步。
\end{enumerate}
第3步中“中间隔着NP或S”这一条件，排除的正是同一小句内的主语，它对应上面所说的宾语代词约束；第4步体现由近及远，广度优先使主语先于宾语被提出，体现主语倾向。每个候选都要通过一致性检查，以及已知的非共指约束，才被接受。

以“Mary met John. She thanked him.”为例，求him的先行项，并设已知Mary用she指称、John用he指称。第1、2步从him的NP上升到当前句的S；第3步在路径左侧遇到She的NP，但它与S之间没有其他NP或S，不提出。当前S已是最高S，第4步转到前一句，广度优先先提出主语Mary，性不匹配被拒绝，随后提出宾语John并接受。输出是him指向John这次提及的指针。若求She的先行项，同样的过程会在前一句先提出Mary并接受。

Hobbs算法不需要训练，结果完全由句法树、搜索顺序和一致性词表决定，这使它便于解释，也让它成为后来学习方法的常用基线和特征来源。代价也来自同一处：句法树错了，候选次序随之改变；两个候选都满足形式条件时，搜索顺序可能选到语义上不通的一个，例如“The trophy didn't fit in the suitcase because it was too big”中，it指trophy还是suitcase要靠常识判断。算法只处理代词，不处理名词短语之间的共指；中文零代词没有可以出发的代词节点，需先检测省略位置。

\subsubsection{提及对与提及排序模型}
\label{sec:nlp-analysis-mention-ranking}

规则方法难以综合字符串匹配、中心词、距离、语法角色、语义类别等大量线索，而且规则之间的优先次序要人工调定。\term{提及对模型}（mention-pair model）把共指改写为监督分类：对每一对提及$(j,i)$（$j$在$i$之前），用特征向量$\bm x_{j,i}$预测两者是否同指，再由这些两两判断构造簇\scite{nlp-analysis-soon2001}。

训练样本的构造决定了模型学到什么。如果把所有提及对都作为样本，负例会远多于正例。Soon等的做法是：对每个有先行项的提及$i$，只取它在同一簇中最近的先行项$j^*$构成正例，再把$j^*$与$i$之间的每个提及与$i$配对作为负例。设$\sigma$为sigmoid函数，$f(\bm x_{j,i})$为分类器打分，以二元交叉熵为目标函数：
\begin{equation}
  L_{\mathrm{pair}}=-\sum_{(j,i)}\Bigl[y_{j,i}\log\sigma\bigl(f(\bm x_{j,i})\bigr)
   +(1-y_{j,i})\log\bigl(1-\sigma(f(\bm x_{j,i}))\bigr)\Bigr],
  \label{eq:nlp-analysis-mention-pair}
\end{equation}
其中$y_{j,i}\in\{0,1\}$表示两者是否同指。Soon等的分类器是C5决策树，特征只有12个，可以分为四组：距离（两者相隔的句子数）；提及形式（$j$、$i$各自是否为代词，$i$是否为有定名词短语或指示性名词短语，两者是否都是专名）；一致性（数、性、语义类别是否相容）；字符串与句法关系（去掉限定词后字符串是否相同，一个是否为另一个的别名或缩写，两者是否构成同位语）。字符串匹配与别名特征承担了专名之间的大部分连接，代词则主要依靠距离与一致性。后来的工作在此基础上增加了中心词、语法角色、句法路径等上百个特征，模型框架保持不变。

推断时，对每个提及$i$从近到远检查前面的提及，第一个概率超过0.5的就作为先行项，称为最近优先聚类；Ng与Cardie改为在所有超过阈值的候选中取概率最高者，称为最佳优先聚类\scite{nlp-analysis-ng-cardie2002}。最后把选中的连接取传递闭包，得到簇。

两两判断互相独立，带来局部决策之间的冲突。设A、B、C三个提及，分类器给出$p(A,B)=0.9$、$p(B,C)=0.8$、$p(A,C)=0.1$。若对所有超过阈值的对取连通分量，A--B与B--C两条边会把三者合成一簇，而模型其实相当确信A与C不同指；传递闭包修复了输出格式，却放大了一次错连。另一个问题是，二分类只问“这对是否同指”，不问“在所有候选中哪个最好”，因此一个代词可能有多个候选都略高于阈值，模型并不比较它们。

\term{提及排序模型}（mention-ranking model）直接回答后一个问题：对提及$i$，把所有较早提及与一个虚拟先行项$\epsilon$放在一起，做一次归一化，选出最合适的一个\scite{nlp-analysis-denis-baldridge2008}。$\epsilon$表示$i$不指向任何较早提及，即开启一个新簇（或$i$不是提及）。记候选集合$Y(i)=\{\epsilon,1,\ldots,i-1\}$，打分函数$s(i,j)$，并约定$s(i,\epsilon)=0$作为基准，则
\begin{equation}
  p(j\mid i)=\frac{\exp s(i,j)}{\sum_{j'\in Y(i)}\exp s(i,j')}.
  \label{eq:nlp-analysis-antecedent-softmax}
\end{equation}

训练时遇到一个新问题：标注给出的是簇，而不是唯一的先行项。“王老师……小王……他”中，“他”的前两个提及都同簇，两个都是正确先行项。若任意指定其中一个为唯一标签，会惩罚同样合理的另一个。Durrett与Klein把正确先行项当作隐变量，最大化正确先行项的概率之和\scite{nlp-analysis-durrett-klein2013}。设$G(i)\subseteq Y(i)$为与$i$同簇的较早提及集合，若$i$是所在簇的第一个提及则$G(i)=\{\epsilon\}$，边缘对数似然目标为
\begin{equation}
  L_{\mathrm{rank}}=-\sum_i\log\sum_{j\in G(i)}p(j\mid i).
  \label{eq:nlp-analysis-coreference-marginal}
\end{equation}
梯度只要求把概率质量从错误候选移到正确候选集合，至于集合内部怎样分配，模型可以自行决定，通常会集中到最容易判断的那个先行项上。

另一种目标是带代价的间隔损失。Wiseman等区分三类错误：本应开新簇却连到某个先行项（误连，false link），本应连接却选了$\epsilon$（误开新簇），以及连错了先行项（错连），分别赋予代价$\alpha_{\mathrm{FL}}$、$\alpha_{\mathrm{FN}}$、$\alpha_{\mathrm{WL}}$，记为$\Delta(i,y)$，正确选择的代价为0\scite{nlp-analysis-wiseman2015}。令$\ell_i=\arg\max_{j\in G(i)}s(i,j)$为当前得分最高的正确先行项，损失为
\begin{equation}
  L_{\mathrm{margin}}=\sum_i\max_{y\in Y(i)}\Delta(i,y)\bigl(1+s(i,y)-s(i,\ell_i)\bigr).
  \label{eq:nlp-analysis-coreference-margin}
\end{equation}
它要求每个错误候选比最好的正确候选至少低1分，代价越大的错误类型受到的惩罚越重；调节三个代价可以在“多连”与“少连”之间权衡，这对后面讨论的评价指标有直接影响，因为MUC与B³对两类错误的敏感程度不同。

提及排序仍是逐提及的局部决策：“他”选了“小王”，并不知道“小王”所在的簇里还有哪些提及，簇层面的性别、数等属性无法汇总。实体级模型让候选单位从单个提及换成已经形成的部分簇，用簇的汇总表示打分，代价是推断顺序影响结果、训练更复杂。后文的高阶推断用另一种方式引入簇信息：迭代地用先行项分布更新提及表示。

\subsubsection{端到端跨度共指与粗到细推断}
\label{sec:nlp-analysis-span-coreference}

前面的模型都假设提及已经由一个独立的检测器给出。检测器漏掉的提及在后续连接中无从补回，检测器多给的伪提及又会成为干扰候选，检测与连接各自训练，彼此不知道对方的需要。Lee等的端到端模型把全部跨度都视为可能的提及，用同一个目标同时学习“哪些跨度是提及”和“提及指向谁”\scite{lee2017coreference}。

\paragraph{跨度表示与打分}
设文档有$n$个词元，编码器给出上下文表示$\bm h_1,\ldots,\bm h_n\in\mathbb R^{d}$。对宽度不超过$L$的每个跨度$i$，记起点与终点为$s_i,e_i$，把跨度简记为$i$，按第~\ref{sec:nlp-analysis-paradigm-span}小节式~\eqref{eq:nlp-analysis-span-representation}构造跨度表示
\[
  \bm g_i=\bigl[\bm h_{s_i};\bm h_{e_i};\hat{\bm h}_i;\phi(i)\bigr],
\]
其中$\hat{\bm h}_i$是跨度内部词元按注意力加权的平均，作用近似于找出中心词；$\phi(i)$是宽度嵌入。在此基础上定义两个分数。提及分数衡量跨度$i$本身像不像一个提及：
\begin{equation}
  s_m(i)=\bm w_m^{\top}\operatorname{FFNN}_m(\bm g_i).
  \label{eq:nlp-analysis-mention-score}
\end{equation}
先行项分数衡量$j$作为$i$的先行项是否合适：
\begin{equation}
  s_a(i,j)=\bm w_a^{\top}\operatorname{FFNN}_a\bigl([\bm g_i;\bm g_j;\bm g_i\circ\bm g_j;\phi(i,j)]\bigr),
  \label{eq:nlp-analysis-antecedent-score}
\end{equation}
其中$\circ$为逐元素乘积，用于直接表达两个跨度在各维上的相似性；$\phi(i,j)$编码两者的距离（按区间分桶后嵌入），以及说话人是否相同、文档体裁等元信息。$\operatorname{FFNN}$是带ReLU的两层前馈网络，$\bm w_m,\bm w_a$把隐藏层映射为标量。三项相加得到配对分数：
\begin{equation}
  s(i,j)=s_m(i)+s_m(j)+s_a(i,j),\qquad s(i,\epsilon)=0.
  \label{eq:nlp-analysis-span-coreference-score}
\end{equation}
按张量形状看，若$\bm h_t\in\mathbb R^{d}$、宽度嵌入$\phi(i)\in\mathbb R^{d_w}$，则$\bm g_i\in\mathbb R^{3d+d_w}$，保留的$M$个跨度堆叠为$\bm G\in\mathbb R^{M\times(3d+d_w)}$，$s_m$对$\bm G$逐行计算得到长度为$M$的向量；先行项打分的输入$[\bm g_i;\bm g_j;\bm g_i\circ\bm g_j;\phi(i,j)]$长度为$3(3d+d_w)+d_\phi$，对每个提及及其$K$个候选计算一次，得到$M\times K$的分数矩阵，再在每行末尾补上$\epsilon$的0分，按行做softmax。

这一分解的意义在于，$i$或$j$只要有一个不像提及，$s_m$就会把配对分数拉低到$\epsilon$的0分以下，模型便选择$\epsilon$。于是“不是提及”和“是提及但开新簇”都由同一个$\epsilon$吸收，提及检测不再需要单独的标签：训练仍用式~\eqref{eq:nlp-analysis-coreference-marginal}的边缘似然，监督只来自共指簇，提及分数作为配对分数的一部分间接得到梯度。

\paragraph{计算量与剪枝}
直接计算的代价是不可接受的。宽度不限时候选跨度有$O(n^2)$个，任意两个跨度都可能构成先行关系，配对数为$O(n^4)$；即使限制宽度$L$，跨度数为$O(nL)$，配对仍有$O(n^2L^2)$个，每对都要跑一次前馈网络。Lee等用两级剪枝把规模降下来，第一级就是第~\ref{sec:nlp-analysis-paradigm-span}小节所说的按一元分数剪枝，$s_m(i)$即该节的一元分数：先按$s_m(i)$只保留得分最高的$\lambda n$个跨度（并去掉与更高分跨度交叉的跨度），再对每个保留的跨度只考虑最近的$K$个先行项。原文取$L=10$、$\lambda=0.4$、$K=250$，配对数降为$O(\lambda nK)$。剪枝发生在训练和推断两个阶段，被剪掉的跨度得不到连接；式~\eqref{eq:nlp-analysis-coreference-marginal}在这种情况下取$G(i)=\{\epsilon\}$，不会强迫模型连到已不在候选中的正确先行项。

图~\ref{fig:nlp-analysis-span-coreference}概括了这一流程。从左到右，候选数量在每一级被压缩，越靠右的打分越昂贵。

\begin{figure*}[htbp]
  \centering
  % 端到端跨度共指的计算流程：跨度枚举、提及打分剪枝、粗粒度先行项筛选与精细打分。
\begin{tikzpicture}[
  x=1mm,y=1mm,figure picture,
  stage/.style={figure node,minimum width=26mm,minimum height=15mm,
    align=center,font=\figurefont\scriptsize},
  size note/.style={font=\figurefont\scriptsize,text=FigureMuted,align=center},
  flow/.style={figure flow}
]
  \node[stage,draw=FigureInput] (enc) at (0,0)
    {编码器\\BiLSTM或SpanBERT\\$\bm H\in\mathbb R^{n\times d}$};
  \node[stage,draw=FigureInput] (span) at (32,0)
    {跨度表示$\bm g_i$\\端点、内部注意力\\与宽度嵌入};
  \node[stage,draw=FigureModel] (ment) at (64,0)
    {提及分数$s_m(i)$\\保留前$\lambda n$个};
  \node[stage,draw=FigureModel] (coarse) at (96,0)
    {粗分$s_c(i,j)$\\每个提及保留前$K$个};
  \node[stage,draw=FigureOutput] (fine) at (128,0)
    {精细分$s_a(i,j)$\\对$Y(i)\cup\{\epsilon\}$\\做softmax};
  \draw[flow] (enc) -- (span);
  \draw[flow] (span) -- (ment);
  \draw[flow] (ment) -- (coarse);
  \draw[flow] (coarse) -- (fine);
  \node[size note] at (32,-13) {约$nL$个候选跨度};
  \node[size note] at (64,-13) {$M=\lambda n$个提及};
  \node[size note] at (96,-13) {$M^2$对只需双线性运算};
  \node[size note] at (128,-13) {$MK$对送入前馈网络};
  \node[size note] at (0,-13) {$n$个词元};
\end{tikzpicture}

  \caption{端到端跨度共指的计算流程。编码器输出逐词元表示，宽度不超过$L$的跨度各得到表示$\bm g_i$；提及分数保留$M=\lambda n$个跨度，粗粒度双线性分数为每个提及保留$K$个候选先行项，只有这$MK$对才进入前馈网络计算精细分数，并与虚拟先行项$\epsilon$一起归一化。下方文字标出每一级的候选规模。}
  \label{fig:nlp-analysis-span-coreference}
\end{figure*}

\paragraph{粗到细推断与高阶推断}
按距离保留最近的$K$个先行项是一种粗糙的筛选：远处的正确先行项会被剪掉，近处大量无关跨度却保留下来。Lee等在后续工作中加入一个廉价的粗粒度分数\cite{nlp-analysis-lee2018higher}：
\begin{equation}
  s_c(i,j)=\bm g_i^{\top}\bm W_c\bm g_j,
  \label{eq:nlp-analysis-coarse-score}
\end{equation}
它对全部$M$个保留跨度两两计算，只是一次$M\times M$的矩阵乘法，不需要逐对运行前馈网络。以$s_m(i)+s_m(j)+s_c(i,j)$为依据为每个$i$选出前$K$个候选（此时$K$可以小得多，如50），再对这些候选计算$s_a$，最终分数为$s_m(i)+s_m(j)+s_c(i,j)+s_a(i,j)$。粗分负责召回，精细分负责区分，二者联合训练。

同一篇工作还引入高阶推断，弥补提及排序缺少簇信息的问题。第$t$轮先按当前表示计算先行项分布$p^{(t)}(j\mid i)$，再用它对候选先行项的表示加权，得到“$i$所在簇的期望表示”，并用门控与原表示融合：
\[
  \begin{aligned}
    \bm a_i^{(t)}&=\sum_{j\in Y(i)}p^{(t)}(j\mid i)\,\bm g_j^{(t)},\qquad
    \bm f_i^{(t)}=\sigma\bigl(\bm W_f[\bm g_i^{(t)};\bm a_i^{(t)}]\bigr),\\
    \bm g_i^{(t+1)}&=\bm f_i^{(t)}\circ\bm g_i^{(t)}+(1-\bm f_i^{(t)})\circ\bm a_i^{(t)},
  \end{aligned}
\]
论文只规定$s(i,\epsilon)=0$，公开实现中把$\epsilon$的表示取为$\bm g_i^{(t)}$本身，即选择$\epsilon$的概率越大，$i$的表示越保持不变。迭代一两轮后，“他”的表示混入了“王老师”的信息，再与后续提及打分时便能利用整个簇的属性。

算法~\ref{alg:nlp-analysis-span-coreference}汇总了推断过程。

\begin{algorithm}[粗到细的端到端跨度共指推断]
  \label{alg:nlp-analysis-span-coreference}
  \begin{algorithmic}[1]
    \Require 文档词元$x_1,\ldots,x_n$；最大宽度$L$，保留比例$\lambda$，候选数$K$，高阶迭代轮数$T$
    \Ensure 共指簇集合
    \State 编码得到$\bm h_1,\ldots,\bm h_n$，为全部宽度不超过$L$的跨度计算$\bm g_i$与$s_m(i)$
    \State 按$s_m$降序选取跨度，跳过与已选跨度交叉者，直到选满$M=\lambda n$个，按位置排序
    \State 计算全部$M\times M$的$s_c(i,j)$，对每个$i$只保留$j<i$中$s_m(i)+s_m(j)+s_c(i,j)$最高的$K$个，记为$Y(i)$
    \For{$t=1,\ldots,T$}
      \State 对$j\in Y(i)$计算$s(i,j)=s_m(i)+s_m(j)+s_c(i,j)+s_a(i,j)$，令$s(i,\epsilon)=0$，得到$p^{(t)}(j\mid i)$
      \State 按高阶推断更新全部$\bm g_i$
    \EndFor
    \State 对每个$i$取$\hat y_i=\arg\max_{j\in Y(i)\cup\{\epsilon\}}s(i,j)$
    \State 把$\hat y_i\neq\epsilon$的指针看作无向边，取连通分量，丢弃只含一个跨度的分量，输出簇
  \end{algorithmic}
\end{algorithm}

\paragraph{算例}
设文本已说明“小王”是“王老师”的别称，依次出现$m_1$“王老师”、$m_2$“小王”、$m_3$“他”，以下分数为教学构造。$m_2$到$m_1$的分数为2，到$\epsilon$为0，选$m_1$。$m_3$到$m_1,m_2,\epsilon$的分数为$2,3,0$，按式~\eqref{eq:nlp-analysis-antecedent-softmax}，概率约为$0.259,0.705,0.035$。$G(m_3)=\{m_1,m_2\}$，边缘概率为$0.965$，损失$-\log0.965\approx0.036$；推断选$m_2$，连通分量为$\{m_1,m_2,m_3\}$。若随后出现“另一位王老师”$m_4$，它到前三者的分数为$-2,-3,-1$，均低于$\epsilon$的0分，于是开启新簇。只比较名字会把$m_4$错并到第一个簇，这里是$s_a$中的上下文表示与修饰语“另一位”压低了分数。

\paragraph{编码器的更替}
原始模型的编码器是以GloVe与Turian词向量、字符卷积为输入的BiLSTM。把编码器换成预训练的BERT后，长文档需要切成不超过512个子词（实验中较短的片段反而更好）的不重叠片段分别编码，片段之间不共享注意力，跨片段的先行项只能依靠跨度表示本身连接\scite{nlp-analysis-joshi2019bert}。SpanBERT在预训练中遮蔽连续跨度并用边界表示预测跨度内容（第~\ref{sec:nlp-classical-spanbert}小节），正好强化了式中$\bm h_{s_i},\bm h_{e_i}$对整个跨度的概括能力；以它为编码器，在SpanBERT原论文的对照中，同一共指模型在OntoNotes英文测试集上的CoNLL F1从使用BERT时的77.1提高到79.6\scite{nlp-classical-joshi2020}。

端到端模型的优势是不需要独立的提及检测器和人工特征，代价是显存随文档长度快速增长，剪枝参数需要权衡召回与计算量；OntoNotes不标注单例，模型因而学会不输出单例，迁移到标注单例的数据集时需要调整。

共指也可以按第~\ref{sec:nlp-analysis-paradigm-generation}小节的约束生成范式求解。Bohnet等把共指写成转移系统的动作序列，再线性化为文本：模型逐句读入文档，对当前句输出一串动作：把某个提及追加到已有的编号簇，或把它与某个较早的提及连接成新簇；动作中的提及以原文字符串加其后的三个词表示，生成后再对齐回原文跨度\scite{nlp-analysis-bohnet2023}。这种做法不需要跨度枚举与剪枝，可以直接利用大型序列到序列模型，代价是解码按句串行、速度较慢，输出的提及字符串必须能在原文中唯一定位。

中文零代词需要在上述框架之外先解决位置问题。常见做法是在谓词前的主语与宾语空位上生成候选零代词位置，判断该位置是否确有省略，再像代词一样从前文的名词短语中选先行项；Chen与Ng用前馈网络分别学习零代词位置与候选先行项的表示并打分\scite{nlp-analysis-chen-ng2016}。位置检测与先行项选择两步的错误会叠加，评价时须说明零代词位置是人工给定还是系统检测的。

\subsubsection{共指评价指标}
\label{sec:nlp-analysis-coreference-metrics}

共指的输出是一组簇，参考答案也是一组簇，比较两组划分没有唯一自然的方式。一个错误的连接可能把两个大实体合并，也可能只多连了一个代词，不同指标对这两类错误给出的惩罚差别很大。设参考簇（key）为$K_1,\ldots,K_a$，系统簇（response）为$R_1,\ldots,R_b$，下面用同一个例子计算三种常用指标。例子有五个提及：$a$“王老师”、$b$“小王”、$c$“他”、$d$“李同学”、$e$“她”。参考答案为$K_1=\{a,b,c\}$、$K_2=\{d,e\}$；系统把“他”错连到李同学，输出$R_1=\{a,b\}$、$R_2=\{c,d,e\}$。

\paragraph{MUC}
MUC指标以连接为单位\scite{nlp-analysis-vilain1995}。一个含$|K|$个提及的簇至少需要$|K|-1$条连接才能连通；系统输出把它切成$|p(K)|$块，其中$p(K)$表示用系统簇划分$K$得到的部分（系统中没有出现的提及各自成为一块），于是系统恢复了$|K|-|p(K)|$条必要连接。召回率与精确率为
\begin{equation}
  R_{\mathrm{MUC}}=\frac{\sum_k\bigl(|K_k|-|p(K_k)|\bigr)}{\sum_k\bigl(|K_k|-1\bigr)},\qquad
  P_{\mathrm{MUC}}=\frac{\sum_l\bigl(|R_l|-|p'(R_l)|\bigr)}{\sum_l\bigl(|R_l|-1\bigr)},
  \label{eq:nlp-analysis-muc}
\end{equation}
其中$p'(R)$是用参考簇划分$R$得到的部分。例子中，$K_1$被系统分成$\{a,b\}$与$\{c\}$两块，恢复$3-2=1$条连接；$K_2$完整地落在$R_2$中，恢复1条，召回率$(1+1)/(2+1)=2/3$。反过来，$R_1$落在$K_1$中，贡献1；$R_2$被参考分成$\{c\}$和$\{d,e\}$，贡献$3-2=1$，精确率$2/3$，F1为$0.667$。

MUC只数连接，带来两个已知缺陷：单例簇不需要任何连接，对分子分母都没有贡献，因此MUC完全忽略单例；把所有提及并成一个大簇只增加少量“多余连接”，精确率下降有限，召回率却达到1，误合并受到的惩罚偏轻。

\paragraph{B³}
B³以提及为单位\scite{nlp-analysis-bagga1998}。对每个提及$m$，记$K(m)$与$R(m)$分别为它所在的参考簇和系统簇，比较两者的重合程度，再对全部$N$个提及取平均：
\begin{equation}
  R_{\mathrm{B^3}}=\frac1N\sum_m\frac{|K(m)\cap R(m)|}{|K(m)|},\qquad
  P_{\mathrm{B^3}}=\frac1N\sum_m\frac{|K(m)\cap R(m)|}{|R(m)|}.
  \label{eq:nlp-analysis-bcubed}
\end{equation}
例子中，召回率各项为：$a,b$各$2/3$，$c$为$1/3$，$d,e$各$2/2=1$，平均$11/15\approx0.733$；精确率各项为：$a,b$各1，$c$为$1/3$，$d,e$各$2/3$，平均同样为$11/15$。B³让每个提及都参与，单例也计入；大簇里的提及多，因而在平均中权重也大。

\paragraph{CEAF}
CEAF先在参考簇与系统簇之间求一个一对一的最优对齐，再按对齐结果计算得分，每个簇最多被使用一次，避免了一个系统簇同时“匹配”多个参考簇\scite{nlp-analysis-luo2005}。设$\phi(K,R)$为一对簇的相似度，$g$为从参考簇到系统簇的一对一映射，$g^*$使总相似度$\Phi(g)=\sum_k\phi(K_k,R_{g(k)})$最大，可用匈牙利算法求得。召回率与精确率为
\begin{equation}
  R_{\mathrm{CEAF}}=\frac{\Phi(g^*)}{\sum_k\phi(K_k,K_k)},\qquad
  P_{\mathrm{CEAF}}=\frac{\Phi(g^*)}{\sum_l\phi(R_l,R_l)}.
  \label{eq:nlp-analysis-ceaf}
\end{equation}
相似度有两种常用取法。基于提及的$\phi_3(K,R)=|K\cap R|$数共有提及，记为CEAF$_m$；基于实体的$\phi_4(K,R)=2|K\cap R|/(|K|+|R|)$对每对簇归一化，记为CEAF$_e$，此时$\phi_4(K,K)=1$，分母就是参考簇数与系统簇数，每个实体同等重要，不论大小。例子中，$\phi_4(K_1,R_1)=2\cdot2/5=0.8$，$\phi_4(K_1,R_2)=2/6\approx0.333$，$\phi_4(K_2,R_1)=0$，$\phi_4(K_2,R_2)=2\cdot2/5=0.8$。最优对齐为$K_1\to R_1$、$K_2\to R_2$，$\Phi=1.6$，召回率与精确率都是$1.6/2=0.8$。

\paragraph{CoNLL F1与对比}
三种指标各有偏向，CoNLL-2011/2012共享任务取MUC、B³、CEAF$_e$三者F1的算术平均作为主指标，称为CoNLL F1\scite{nlp-analysis-pradhan2012}。例子的CoNLL F1为$(0.667+0.733+0.8)/3\approx0.733$。不同实现对未匹配到参考的系统提及处理不一，Pradhan等给出了参考实现，报告结果时应使用同一评分脚本\scite{nlp-analysis-pradhan2014}。

表~\ref{tab:nlp-analysis-coreference-metrics}把同一参考答案下的三种系统输出放在一起：上面的错连输出、把五个提及全部并成一簇的误合并输出，以及每个提及各自成簇的全单例输出。

\begin{table}[htbp]
  \centering
  \small
  \begin{tabular}{@{}lcccc@{}}
    \toprule
    系统输出 & MUC F1 & B³ F1 & CEAF$_e$ F1 & CoNLL F1\\
    \midrule
    错连：$\{a,b\},\{c,d,e\}$ & 0.667 & 0.733 & 0.800 & 0.733\\
    全部合并：$\{a,b,c,d,e\}$ & 0.857 & 0.684 & 0.500 & 0.680\\
    全单例：$\{a\},\{b\},\{c\},\{d\},\{e\}$ & 0 & 0.571 & 0.333 & 0.302\\
    \bottomrule
  \end{tabular}
  \caption{同一参考答案$\{a,b,c\},\{d,e\}$下三种系统输出的共指指标（教学算例）。}
  \label{tab:nlp-analysis-coreference-metrics}
\end{table}

表中最值得注意的是第二行：全部合并在MUC上的得分高于只错连一个代词的系统，因为它召回了全部连接，只多出一条；B³与CEAF$_e$则明显惩罚这种误合并，CEAF$_e$因为只能把一个系统簇对齐到一个参考簇，惩罚最重。第三行相反，全单例在MUC上为0，却在B³上拿到0.571，因为每个提及至少与自己重合。没有哪一个指标单独可靠，这正是取三者平均的理由；分析错误时仍应分别查看三个指标，以及提及检测本身的召回率与精确率。Moosavi与Strube指出，B³与CEAF在系统提及与参考提及不一致时也会出现反直觉的结果，并提出按实体大小加权、以连接为单位计分的LEA指标\cite{nlp-analysis-moosavi2016}；目前英文共指的主要报告口径仍是CoNLL F1。

这些指标最常用的数据是CoNLL-2012共享任务的OntoNotes~5.0语料，它包含英语、汉语和阿拉伯语三部分，英语部分约有2800篇训练文档、约340篇开发文档与约350篇测试文档，体裁覆盖新闻、广播对话、电话交谈、网络文本等。OntoNotes不标注单例，也不标注与通指名词短语的共指，因此在它上面训练的系统默认不输出单例。

表~\ref{tab:nlp-analysis-coreference-methods}对照本小节的四类方法。

\begin{table*}[htbp]
  \centering
  \small
  \begin{tabularx}{\linewidth}{@{}l XXX@{}}
    \toprule
    方法 & 解决的问题 & 核心机制 & 新的代价或局限\\
    \midrule
    Hobbs算法 & 无标注数据时消解代词 & 句法树上的固定搜索顺序加一致性筛选 & 依赖句法树质量；只处理代词；不用语义\\
    提及对模型 & 综合大量特征 & 两两二分类，最近优先或最佳优先聚类 & 判断互相独立，传递闭包放大错连；候选间不比较\\
    提及排序模型 & 候选间的比较 & 对全部候选与$\epsilon$做softmax，正确先行项作隐变量 & 仍是逐提及决策，缺少簇级信息\\
    端到端跨度模型 & 提及检测与连接的流水线误差 & 跨度表示、$s_m+s_a$分解、剪枝与粗到细筛选 & 显存随长度增长；剪枝损失召回\\
    \bottomrule
  \end{tabularx}
  \caption{共指消解方法对照。从上到下，每一行解决上一行留下的主要问题。}
  \label{tab:nlp-analysis-coreference-methods}
\end{table*}

\subsection{篇章单元与关系}
\label{sec:nlp-analysis-discourse-relations}

身份共指连接对象，篇章关系连接承载内容的单元。\term{基本篇章单元}（elementary discourse unit, EDU）是标注规范选定的最小分析片段，通常接近一个分句。用“下雨了，所以路面湿了。车辆减速。”作为本小节的统一材料，切成三个单元：$u_1$“下雨了，”、$u_2$“所以路面湿了。”、$u_3$“车辆减速。”。$u_1$与$u_2$之间有连接词“所以”，表达结果关系；$u_2$与$u_3$之间没有连接词，读者仍会理解为结果关系。这里分析的是文字所表达的关系，不是对现实因果的验证。

篇章关系有两种主要的标注框架。宾州篇章树库（Penn Discourse Treebank, PDTB）以连接词为中心，标注一个连接词及其两个论元（Arg1、Arg2）和关系类别，不要求全文形成一棵树\scite{nlp-analysis-prasad2008}。PDTB的关系类别分三级：一级为时序（Temporal）、偶然（Contingency）、比较（Comparison）、扩展（Expansion）四类，二级细分为十余种类型（隐式关系分类常用其中11类），三级再区分方向等细节。没有显式连接词的相邻句对，标注者会补写一个最合适的连接词，再按它标注关系，这类关系称为\term{隐式篇章关系}（implicit discourse relation）。\term{修辞结构理论}（rhetorical structure theory, RST）则要求用关系把全部EDU递归地组合成覆盖全文的树，每个内部节点记录关系，以及子节点的\term{核性}（nuclearity）：核心（nucleus）承载作者的主要意图，附属（satellite）为核心提供原因、背景、阐述等支持\scite{nlp-analysis-mann-thompson1988}。RST篇章树库（RST-DT）在《华尔街日报》的385篇文章上按这一理论标注，78种细粒度关系被归为16类\scite{nlp-analysis-carlson2001}；解析研究通常再加上文本组织（Textual-Organization）与同一单元（Same-Unit）两种结构关系，按18类评价。

篇章关系分析的挑战可以列为五项：单元切分，即EDU边界在哪里；连接词识别与消歧，同一个词可能不在篇章层面起连接作用，起连接作用时也可能表达多种关系（英语since既可表原因也可表时间）；隐式关系，缺少连接词这一最强信号；树结构，RST要求层次决策互相一致；标注稀缺，RST-DT只有数百篇文档。方法谱系与之对应：边界标注加显式连接词分类处理有信号的情形；隐式关系分类用论元对编码器从两段的语义推断关系；RST移进--归约按第~\ref{sec:nlp-analysis-paradigm-transition}小节的转移范式自底向上构建整棵树；递归分割则自顶向下决定树的划分，避免早期局部合并的错误一直保留到最后。

\subsubsection{篇章单元切分与显式关系分类}
\label{sec:nlp-analysis-discourse-boundary}

EDU切分可以写成序列标注：每个字或词标B（单元开始）或I（单元内部），按第~\ref{sec:nlp-analysis-paradigm-tagging}小节的方法训练，句首位置强制为B。上面16个字符的材料在位置0、4、11取B，其余取I，解码后恢复$[0,4)$、$[4,11)$、$[11,16)$三个单元。标点在本例归入左侧单元，不同规范可以不同，但恢复出的单元必须覆盖全部原文，且彼此不重叠。英语RST-DT上的EDU切分主要依赖句法与标点，用预训练编码器加序列标注已经相当准确，篇章解析中的主要误差来自后面的结构与关系决策。

显式关系分类分两步。第一步判断候选连接词是否真在篇章层面起作用：“所以”多为连接用法，“而且”“然后”却也常出现在小句内部；英语and有大量非篇章用法。Pitler与Nenkova用连接词本身及其句法上下文（所在成分的类别、父节点和左右兄弟节点）作为特征，区分篇章用法与非篇章用法，并在此基础上判断关系类别\scite{nlp-analysis-pitler-nenkova2009}。第二步给已确认的连接词分配关系类别：大多数连接词的关系分布高度集中，在PDTB上仅用连接词本身作特征，一级四分类的准确率就超过九成\scite{nlp-analysis-pitler2008}。这一事实说明，显式关系的难点不在关系分类本身，而在论元跨度的定位，尤其是Arg1可能在前几句之中。

对$u_1,u_2$，编码两段及连接词，在原因、对比、阐述等关系上做softmax。设三个logit为$2,0,-1$，原因关系的概率约为$0.844$，输出“原因论元$u_1$，结果论元$u_2$，连接词‘所以’”的记录。关系的方向必须与论元对应：把$u_2$当原因论元，类别名正确也是错误的输出。因此评价时应分别报告连接词识别、论元跨度和关系三项。

把这些步骤串起来，就得到端到端的PDTB式篇章解析器。Lin等的流水线依次执行：连接词分类器判断候选词是否为篇章用法；论元位置分类器判断Arg1与连接词在同一句还是在前面的句子中；论元抽取器在句法树上选出两个论元的跨度；显式关系分类器给出类别；最后对相邻句对中没有显式连接词的位置，再判断是否存在隐式关系及其类别\scite{nlp-analysis-lin2014pdtb}。流水线每一步都依赖上一步的输出，端到端结果明显低于各子任务在给定前一步正确输入时的成绩，这也是DISRPT等评测把分割、连接词识别与关系分类分别作为独立任务报告的原因\scite{braud-etal-2023-disrpt}。

\subsubsection{隐式篇章关系分类}
\label{sec:nlp-analysis-implicit-relation}

隐式关系难在什么地方，可以对照显式关系看。显式关系中，连接词几乎决定了类别；隐式关系中，这个信号不存在，模型只能从两段的内容推断联系。“路面湿了”与“车辆减速”之间的结果关系，依赖“湿滑路面导致减速”这一常识；把$u_3$换成“行人打着伞”，关系就变为并列或阐述。推断需要比较两段的语义，而不是只看其中一段。标注本身也更难，标注者在补写连接词时意见分歧较大。因此在PDTB上，隐式关系分类的准确率长期远低于显式关系\scite{nlp-analysis-pitler2009}。

早期方法用词对特征：把Arg1中的每个词与Arg2中的每个词组合成特征，期望“湿—减速”这类词对反映关系。词对数量巨大且稀疏，泛化很差。神经方法改用\term{论元对编码器}（argument pair encoder）：分别编码两个论元，再让它们互相注意，最后分类。下面给出一个代表性结构的逐层运算。

设两个论元的词嵌入为$\bm X^{(1)}\in\mathbb R^{n_1\times d}$与$\bm X^{(2)}\in\mathbb R^{n_2\times d}$。第一层用共享参数的BiLSTM分别编码，得到$\bm H^{(1)}\in\mathbb R^{n_1\times 2h}$与$\bm H^{(2)}\in\mathbb R^{n_2\times 2h}$。第二层计算交互矩阵并做双向软对齐：
\begin{equation}
  \begin{aligned}
    \bm E&=\bm H^{(1)}\bm W\bm H^{(2)\top}\in\mathbb R^{n_1\times n_2},\\
    \tilde{\bm H}^{(1)}&=\operatorname{softmax}_{\mathrm{row}}(\bm E)\,\bm H^{(2)},\qquad
    \tilde{\bm H}^{(2)}=\operatorname{softmax}_{\mathrm{col}}(\bm E)^{\top}\bm H^{(1)},
  \end{aligned}
  \label{eq:nlp-analysis-implicit-interaction}
\end{equation}
其中$\bm E_{p,q}$衡量Arg1第$p$个词与Arg2第$q$个词的关联，$\tilde{\bm H}^{(1)}\in\mathbb R^{n_1\times 2h}$的第$p$行是Arg2中与该词对齐的内容。第三层把原表示与对齐表示组合为$[\bm H^{(1)};\tilde{\bm H}^{(1)};\bm H^{(1)}-\tilde{\bm H}^{(1)};\bm H^{(1)}\circ\tilde{\bm H}^{(1)}]\in\mathbb R^{n_1\times 8h}$，差与积直接表达两侧的差异与一致；对它沿词的维度做最大池化与平均池化，得到$\bm r^{(1)}\in\mathbb R^{16h}$，Arg2同理得到$\bm r^{(2)}$。最后一层把$[\bm r^{(1)};\bm r^{(2)}]$送入多层感知机，在$C$个关系上做softmax，以交叉熵训练。

用一个二维小例子看交互层做了什么（数值为教学构造）。设Arg1为“路面／湿”，Arg2为“车辆／减速”，$\bm H^{(1)}$的两行为$(1,0)$与$(0,1)$，$\bm H^{(2)}$的两行为$(1,0)$与$(0,2)$，取$\bm W=\bm I$。则$\bm E=\bigl(\begin{smallmatrix}1&0\\0&2\end{smallmatrix}\bigr)$，按行softmax得到“路面”对Arg2两词的权重约为$(0.731,0.269)$，“湿”约为$(0.119,0.881)$；对应的对齐表示为$\tilde{\bm h}^{(1)}_1\approx(0.731,0.538)$、$\tilde{\bm h}^{(1)}_2\approx(0.119,1.762)$。“湿”主要对齐到“减速”，后续的差与积特征就围绕这一对词计算，分类器得以利用“湿—减速”这类跨论元的联系，而不必在稀疏的词对特征中查找。

这一结构与第~\ref{sec:nlp-dialogue-decomposable}小节判断证据关系的可分解注意力同属“对齐—比较—汇聚”的句对模型；Liu与Li的多层注意力模型则让两个论元反复互相读取多轮\scite{nlp-analysis-liu-li2016}。

预训练编码器把这一过程交给自注意力。输入“[CLS] Arg1 [SEP] Arg2 [SEP]”，两个论元的词元在每一层都互相注意，取$\bm h_{\mathrm{[CLS]}}$接线性层分类即可。Shi与Demberg在《华尔街日报》文本上继续预训练BERT时比较了保留与去掉下一句预测（NSP）目标两种做法，去掉NSP后隐式关系分类的准确率下降，他们据此认为NSP有助于判断“两段能否前后衔接”\scite{nlp-analysis-shi-demberg2019}。这一比较针对的是继续预训练阶段，原始BERT本身仍经过NSP预训练。

隐式关系分类的评价设置需要说明清楚，否则数字不可比。PDTB 2.0上有两种常用划分：一种以第2--20节为训练集、第0--1节为开发集、第21--22节为测试集，另一种以第2--21节训练、第22节开发、第23节测试；任务可以是一级四类的单一多分类，也可以是对每个一级类别做“一类对其余”的二分类，后者报告各类F1；二级任务通常在11个类型上做多分类并报告准确率。由于类别极不均衡（扩展类占多数），四分类常同时报告宏平均F1。不同论文还在是否包含实体关系（EntRel）、是否把多标签样本拆开计算上有差别，比较前应核对这些设置。

标注数据不足是另一个瓶颈。一个自然的想法是利用大量带显式连接词的文本：去掉连接词，把它当作隐式样本，标签由连接词决定\scite{nlp-analysis-marcu-echihabi2002}。Sporleder与Lascarides的实验说明这样得到的样本与真实隐式样本分布不同：作者省略连接词的句对，往往是关系本来就容易从内容看出的句对，而人为删去连接词的句对，内容未必支持该关系，于是在人造样本上训练的分类器迁移到真实隐式关系时效果很差\scite{nlp-analysis-sporleder2008}。较稳妥的用法是把连接词作为中间预测目标：先让模型预测“这里最适合补写哪个连接词”，再由连接词推断关系，或者把连接词预测作为辅助任务与关系分类联合训练\scite{nlp-analysis-zhou2010}。

\subsubsection{RST移进--归约篇章解析}
\label{sec:nlp-analysis-rst-transition}

RST要求把全部EDU组合成一棵树。自底向上的移进--归约是RST解析的主要做法之一，Ji与Eisenstein把它与可学习的EDU表示结合，在每个状态上用学到的表示选择动作\scite{nlp-analysis-ji-eisenstein2014}。它与第~\ref{sec:nlp-analysis-shift-reduce}小节的移进--归约成分解析在结构上同构：叶节点由词换成EDU，短语类别换成“关系＋核性”，状态仍由子树栈和待读队列组成。因此这里只写两者的差异。

动作集合为SHIFT与REDUCE$(r,\nu)$。SHIFT把队首EDU入栈；REDUCE$(r,\nu)$弹出栈顶两棵相邻子树，以关系$r$和核性$\nu\in\{\mathrm{NS},\mathrm{SN},\mathrm{NN}\}$组成父树后入栈，其中NN对应列举、对比等多核心关系。与成分解析相比，它没有一元动作，篇章树中不存在单子节点的内部节点；归约动作的标签空间是关系数乘以核性组合，约为$18\times3$量级，远大于成分解析的短语类别数，而训练样本少得多。为此，Wang等提出两阶段方法：第一阶段的转移系统只预测不带关系的树结构与核性，动作只有SHIFT与三种REDUCE；第二阶段再为每个内部节点分类关系，并对句内、跨句、跨段落的节点使用不同的分类器\scite{nlp-analysis-wang2017twostage}。

表~\ref{tab:nlp-analysis-rst-actions}列出两种系统动作的对应关系。

\begin{table}[htbp]
  \centering
  \small
  \begin{tabularx}{\linewidth}{@{}l XX@{}}
    \toprule
    & 成分移进--归约 & RST移进--归约\\
    \midrule
    叶节点 & 词 & EDU\\
    SHIFT & 队首词入栈 & 队首EDU入栈\\
    归约 & REDUCE-L/R-$X$：合并栈顶两项，标短语类别$X$，记录中心词方向 & REDUCE$(r,\nu)$：合并栈顶两项，标关系$r$与核性$\nu$\\
    一元动作 & UNARY-$X$ & 无\\
    子树代表 & 中心词 & 核心EDU\\
    终止 & 队列空且栈中一棵树 & 同左\\
    \bottomrule
  \end{tabularx}
  \caption{成分移进--归约与RST移进--归约的动作对照。核性中的N方向对应成分解析中的中心词方向。}
  \label{tab:nlp-analysis-rst-actions}
\end{table}

训练与解码沿用第~\ref{sec:nlp-analysis-paradigm-transition}小节：参考树经后序遍历给出唯一的静态oracle动作序列，在每个状态上对合法动作做softmax，以式~\eqref{eq:nlp-analysis-transition-loss}的交叉熵训练；栈中少于两棵子树时REDUCE非法，队列为空时SHIFT非法。状态特征取栈顶两棵子树与队首EDU的表示，子树表示可由其核心EDU的表示代表，这与成分解析中用中心词代表短语是同一种思路；此外还加入两棵子树是否在同一句、同一段落等结构特征，因为句内关系与跨段落关系的分布差别很大。

对三个单元，参考树为$[[u_1^{\mathrm S}u_2^{\mathrm N}]_{\text{原因}}^{\mathrm S}u_3^{\mathrm N}]_{\text{原因}}$，动作序列为SHIFT、SHIFT、REDUCE(原因,SN)、SHIFT、REDUCE(原因,SN)。另一条合法轨迹SHIFT、SHIFT、SHIFT、REDUCE、REDUCE先合并$u_2,u_3$，得到右分支的树。设前一条轨迹的动作分数之和为8，后一条为6，束搜索在完整轨迹之间比较时会选前者；贪心解码若在第三步就因局部分数选择了SHIFT，此后再无法组成$[u_1u_2]$这一子树。

\paragraph{RST-Parseval评价}
RST树的评价沿用成分句法的Parseval思想，比较参考树与预测树中的跨度。Marcu提出的RST-Parseval把每个非根节点（包括叶节点EDU）看作一个评价单位，并给它附上核性与关系：核性是该节点在其父节点中是核心还是附属，关系是它与兄弟节点之间的关系\scite{nlp-analysis-marcu2000}。在此基础上逐级加严，得到四个指标：只比较跨度（S），跨度与核性都对（N），跨度与关系都对（R），三者都对（F）。EDU切分给定时，两棵树的节点数相同，精确率等于召回率。

用三个单元的例子计算。参考树的四个非根节点为$u_1$（S，原因）、$u_2$（N，原因）、$[u_1u_2]$（S，原因）、$u_3$（N，原因）。设预测树为$[u_1^{\mathrm S}[u_2^{\mathrm N}u_3^{\mathrm S}]_{\text{阐述}}^{\mathrm N}]_{\text{原因}}$，其非根节点为$u_1$（S，原因）、$[u_2u_3]$（N，原因）、$u_2$（N，阐述）、$u_3$（S，阐述）。跨度上，三个叶节点都匹配，$[u_1u_2]$与$[u_2u_3]$不匹配，S为$3/4$；核性上$u_1,u_2$匹配，N为$2/4$；关系上只有$u_1$匹配，R为$1/4$；F同样为$1/4$。预测树唯一的内部结构决策是错的，跨度指标却有0.75，原因是叶节点在EDU给定时总是匹配。Morey等据此指出，RST-Parseval因计入叶节点而抬高了分数，且不同论文混用不同的计算方式，使结果不可直接比较；他们建议改用只评价内部节点的标准Parseval，并统一评分脚本\scite{nlp-analysis-morey2017}。

\subsubsection{递归分割式篇章树恢复}
\label{sec:nlp-analysis-rst-split}

移进--归约在看到全文之前就要作出局部合并，较早的错误合并会一直保留在树中。自顶向下的方法反过来，先在整段上决定最高层的划分，再递归处理两侧。最高层的划分往往对应段落或语义大块之间的边界，可以利用全文信息作出判断\scite{nlp-analysis-lin2019unified,nlp-analysis-zhang2020topdown,nlp-analysis-kobayashi2020}。

设EDU边界区间$[i,j)$包含$j-i$个单元，若$j-i=1$直接返回叶节点；否则在$i<k<j$中选择分割点，再预测核性$\nu$与关系$r$，递归处理$[i,k)$与$[k,j)$。实现上，编码器先给出每个EDU的表示$\bm e_1,\ldots,\bm e_m$，区间表示$\bm z_{i,j}$可由区间内EDU表示汇聚得到（或用解码器的隐状态逐步维护）；分割头对候选边界做指针注意力：
\[
  p(k\mid i,j)=\operatorname{softmax}_{i<k<j}\bigl(\bm v^{\top}\tanh(\bm W_1\bm z_{i,j}+\bm W_2\bm e_k)\bigr),
\]
核性头与关系头在选定的分割下，以两个子区间的表示为输入分类。参考树的每个内部节点提供分割、核性、关系三项监督，各项交叉熵相加作为目标函数。区间按栈顺序依次展开，$m$个EDU共需$m-1$次分割决策。

对$[0,3)$，设分割点$k=1,2$的logit为1与3，选$k=2$，概率约$0.881$，核性SN与关系“原因”得分最高；子区间$[0,2)$只有$k=1$一种分割，仍取原因与SN；$[2,3)$等单元区间直接返回。恢复的树与第~\ref{sec:nlp-analysis-rst-transition}小节的参考树相同。若最高层误选了$k=1$，那么$[0,2)$这一跨度不可能在后续出现，这说明自顶向下方法同样有错误传播，只是方向相反：高层错误限制了所有下层结构。

如果分割分数只依赖当前区间与分割点，不依赖此前的分割历史，那么整棵树的分数是各内部节点分数之和，可以像第~\ref{sec:nlp-analysis-neural-cky}小节的神经CKY（式~\eqref{eq:nlp-analysis-neural-cky}）那样，以EDU为单位用区间动态规划求精确最优：
\[
  D(i,j)=\max_{i<k<j,\,r,\nu}\bigl\{s(i,j,k,r,\nu)+D(i,k)+D(k,j)\bigr\},\qquad D(i,i+1)=0,
\]
保存每个区间的最优分割与标签即可回溯，复杂度为$O(m^3)$。贪心分割只需$O(m)$次决策，但不保证全局最优；解码器若用隐状态记录了此前的分割历史，分数不再可分解，上式就不能直接使用。

表~\ref{tab:nlp-analysis-discourse-methods}对照本小节的方法。

\begin{table*}[htbp]
  \centering
  \small
  \begin{tabularx}{\linewidth}{@{}l XXX@{}}
    \toprule
    方法 & 解决的问题 & 核心机制 & 新的代价或局限\\
    \midrule
    边界标注与显式关系分类 & EDU切分；连接词消歧与关系分类 & B/I标注；连接词及其句法上下文特征 & 论元跨度定位仍难，尤其Arg1远离连接词时\\
    论元对编码器 & 无连接词的隐式关系 & 两论元交互注意力或句对编码，再分类 & 依赖常识；以显式样本扩充时分布不同\\
    RST移进--归约 & 覆盖全文的层次树 & 转移系统，归约动作带关系与核性 & 局部早期合并不可撤回；标签空间大\\
    递归分割 & 自底向上的局部合并错误 & 自顶向下指针选择分割点；可分解时用区间DP & 高层分割错误限制下层；精确解需分数无历史依赖\\
    \bottomrule
  \end{tabularx}
  \caption{篇章单元与关系分析方法对照。}
  \label{tab:nlp-analysis-discourse-methods}
\end{table*}

\subsection{论点、理由与论证结构}
\label{sec:nlp-analysis-argumentation}

篇章关系中的“解释”说明两段内容如何衔接，论证分析关心的则是一段文字是否被用来支持或反驳另一个主张。\term{论证分析}（argument mining）先定位论证单元并判断类型，再恢复单元之间有方向的支持与攻击关系。以下材料含三个单元：$a_1$“应增加自习座位”，$a_2$“晚间常有人找不到空位”，$a_3$“本月晚间到场人数持续上升”。在标注中，$a_1$是主张，$a_2$支持$a_1$，$a_3$支持$a_2$。论证结构描述作者的推理组织，不判断这些陈述是否属实，也不判断理由是否充分。

这一小节的数据主要来自Stab与Gurevych标注的议论文语料（第2版含402篇学生议论文），单元类型为主论点（MajorClaim）、论点（Claim）和前提（Premise），关系为支持与攻击，论点相对主论点另标赞成或反对的立场。

论证结构分析的挑战有三项。单元边界与类型：论证单元不一定与句子或EDU对齐，同一句可能包含主张与理由；类型（主论点、论点、前提）依赖它在整篇中的位置。关系方向：支持关系从理由指向主张，方向错误等于结构错误。全局结构：在Stab与Gurevych的议论文语料中，每段的论证结构是以主张为根的树或森林，局部预测的关系可能成环或给一个前提指向多个目标\scite{stab-gurevych-2017-parsing}。方法谱系有两条：流水线先分类单元与关系，再用整数线性规划（ILP）恢复全局合法的结构；双仿射依存解析把单元当作依存树的节点，用神经弧评分接同一套结构约束，并可与单元识别联合。

\subsubsection{论证单元分类与ILP结构恢复}
\label{sec:nlp-analysis-argument-ilp}

Stab与Gurevych的流水线分四步：用序列标注识别论证单元，分类单元类型，对单元对预测是否存在连接，最后用ILP联合协调前两类局部输出。ILP中边的权重综合了关系分类器的连接概率、由类型分类器给出的“该单元像主张”的程度，以及由其他边推出的入边与出边数。求解后，没有出边的单元为主张，其余为前提；支持或攻击的立场再单独分类\scite{stab-gurevych-2017-parsing}。

用二元变量$x_{i,j}\in\{0,1\}$表示是否选择从理由$i$到目标$j$的边，$w_{i,j}$为边的权重，目标是最大化所选边的总权重：
\[
  \max_{\bm x}\sum_{i\neq j}w_{i,j}x_{i,j}.
\]
基本约束为
\[
  x_{i,i}=0,\qquad \sum_jx_{i,j}\leq1\ \ (\forall i),\qquad \sum_{i,j}x_{i,j}\leq n-1,
\]
分别表示没有自环、每个单元至多指向一个目标、总边数不超过树的边数。这三项还不能排除环，例如$a_2\to a_3$与$a_3\to a_2$同时满足它们。为此引入层次变量$0\leq u_i\leq n-1$，要求
\[
  u_i\geq u_j+1-n(1-x_{i,j}),\qquad i\neq j.
\]
选中边$i\to j$时约束变为$u_i\geq u_j+1$，即理由的层次严格高于目标；沿一个环累加这些不等式会得到$u_i\geq u_i+k$（$k$为环长），矛盾，因此不能成环。未选中时约束右端至多为$u_j+1-n\leq0$，自动满足。这样得到的是有向森林；若标注要求每段恰好一棵树，再加上总边数等于$n-1$即可。

对三个单元，设主张$a_1$不能有出边，其余候选边的权重为$w_{2,3}=6$、$w_{3,2}=5$、$w_{2,1}=4$、$w_{3,1}=2$（教学构造）。各自取最高分的出边，得到$a_2\to a_3$与$a_3\to a_2$，总分11，却成环。合法组合中，$a_2\to a_1,a_3\to a_2$得9分，$a_2\to a_3,a_3\to a_1$得8分，ILP取9分的结构：$a_1$为根，$a_2,a_3$为前提。

ILP的优点是约束可以按标注规范自由增减，局部分类器可以任意替换；代价是求解时间在单元数多时上升，而且流水线中单元识别的错误无法由ILP修复。两个理由共同才能支持一个结论（合取前提）的情形，普通二元边无法表达，需要额外的论证节点。

\subsubsection{双仿射论证依存解析}
\label{sec:nlp-analysis-argument-biaffine}

论证结构是以主张为根、每个前提至多指向一个目标的树或森林，这与依存树“每个词至多一个中心词”的约束相同。因此可以把论证单元当作词，把“前提指向目标”当作“依存词指向中心词”，直接复用第~\ref{sec:nlp-analysis-biaffine}小节的双仿射依存解析\scite{nlp-analysis-eger2017}。每个单元先编码为向量$\bm u_i$（如对单元内词元表示做池化），再经两个多层感知机分别投影为“依据端”表示$\bm h_i^{(d)}$与“目标端”表示$\bm h_j^{(h)}$，弧分数按式~\eqref{eq:nlp-analysis-biaffine}计算为$s_{i,j}=\bm h_j^{(h)\top}\bm U\bm h_i^{(d)}+\bm h_j^{(h)\top}\bm u$。类型头预测主张、前提或非论证，关系头预测支持或攻击。训练对标注的弧、类型与关系分别计算交叉熵；解码时把$s_{i,j}$作为第~\ref{sec:nlp-analysis-argument-ilp}小节整数规划中的边权$w_{i,j}$，施加同一套森林约束，也可以用Chu--Liu/Edmonds算法求最大生成树。

令依据端$\bm h_2^{(d)}=(1,0)^{\top}$、$\bm h_3^{(d)}=(0,1)^{\top}$，目标端$\bm h_1^{(h)}=(4,2)^{\top}$、$\bm h_2^{(h)}=(0,5)^{\top}$、$\bm h_3^{(h)}=(6,0)^{\top}$，取$\bm U$为单位矩阵、$\bm u=\bm 0$，则$s_{2,1}=4$、$s_{3,1}=2$、$s_{3,2}=5$、$s_{2,3}=6$，恰好是该小节算例的边权，约束解码得到同一个9分结构。双仿射评分替换的是边权的来源，结构约束与解码问题不变。

与ILP流水线相比，双仿射解析让单元表示、弧评分和类型分类共享一个编码器联合训练，不必手工设计单元对特征。单元边界未知时，可以把词元级BIO标注与单元间的弧预测放在同一模型中联合学习\cite{nlp-analysis-eger2017}；Eger等比较了把端到端论证分析写成依存解析、序列标注与多任务学习等几种形式，发现在他们的议论文数据上，直接套用依存解析的效果较差，以BiLSTM为基础、把单元类型与关系距离编进标签的序列标注模型在不同设置下表现较稳健，多任务学习也有帮助；Ye与Teufel进一步把论证单元内部用带类型的连接串起来，以单元的末词作为对外端点，在词元层面做依存解析，再把端点间的弧映射回单元间的关系，这一做法允许非树的论证图\scite{ye-teufel-2021-end}。对于本身就是图（一个前提支持多个主张）的语料，Morio等去掉单父约束，对每对单元独立判断是否存在边\scite{nlp-analysis-morio2020}。约束放松后可以覆盖更多结构，但不再能用最大生成树保证输出合法，需要另外检查环。

\subsection{内容衔接与连贯性}
\label{sec:nlp-analysis-coherence}

前三个小节输出具体的边；连贯性分析转而评价整段文字：给段落一个连贯性分数，或者把打乱的句子恢复成最合理的顺序。比较三句材料：A“小王借了一本书”，B“他读完了这本书”，C“这本书被归还”。原顺序ABC读来自然；顺序ACB让归还先于阅读，所有实体词都还在，读者却感到不通。

连贯性的挑战首先是缺少直接标注：很少有语料给段落标连贯分数。常见做法是用原文作为正例，用打乱句序、删去指代、插入无关句等扰动构造负例，训练模型区分两者。这种监督有两个隐患：模型可能只学会识别某一种扰动，而且可能依赖句首编号、段落长度、标点等与连贯无关的捷径。方法谱系为：中心理论描述相邻句之间注意焦点如何转移；实体网格把这种转移变成可学习的特征；句对衔接评分加束搜索直接求句子排列；指针网络以已选句子的全部历史为条件逐句选择。句子排序的结果用Kendall τ与完全匹配率评价。

\subsubsection{中心理论与实体网格}
\label{sec:nlp-analysis-entity-grid}

\term{中心理论}（centering theory）描述局部连贯：相邻两句如果围绕同一个实体展开，并且该实体保持在显著的语法位置，读者处理起来最省力\scite{nlp-analysis-grosz1995}。对每个话语单元$U_n$（这里取句子）定义三个对象。前瞻中心列表$C_f(U_n)$是$U_n$中出现的实体，按显著性排序，通常主语高于宾语，宾语高于其他成分。回指中心$C_b(U_n)$是$C_f(U_{n-1})$中在$U_n$里再次出现的、排位最高的实体，即这句话“接着谈论”的对象。首选中心$C_p(U_n)$是$C_f(U_n)$的第一个元素，预示下一句最可能继续谈论的对象。

相邻句之间的过渡按两个问题分类：回指中心是否与上一句相同，回指中心是否就是本句的首选中心。两者都是为延续（Continue），前者是而后者否为保持（Retain），前者否为转移（Shift）。中心理论还给出一条关于代词的规则：若$C_f(U_{n-1})$中有实体在$U_n$里用代词实现，则$C_b(U_n)$也必须用代词实现。它解释了为什么“小王借了一本书。小王读完了它”显得别扭：书用了代词，排位更高、作为回指中心的小王却用了全名。Brennan等又把转移细分为$C_b=C_p$与$C_b\neq C_p$两种，并提出偏好次序：延续优于保持，保持优于前一种转移，前一种转移优于后一种\scite{nlp-analysis-brennan1987}；后来的文献把这两种转移分别称为平滑转移（smooth shift）与粗糙转移（rough shift）。

对ABC，$C_f(A)=[\text{小王},\text{书}]$；B中两者都再现，$C_b(B)=\text{小王}=C_p(B)$，A没有回指中心，按惯例记为延续；C中只出现书，$C_b(C)=\text{书}\neq C_b(B)$，且$C_p(C)=\text{书}$，为平滑转移。对ACB，A到C为延续，C到B时$C_b(B)=\text{书}=C_b(C)$，但$C_p(B)=\text{小王}$，为保持。按偏好次序，ACB的过渡序列（延续、保持）反而优于ABC（延续、平滑转移）。这个例子说明中心理论只刻画实体焦点的延续，不刻画事件的时间顺序；ABC优于ACB的理由在于“先读后还”，这属于事件知识。

\term{实体网格}（entity grid）把中心理论的思想变成统计特征\scite{barzilay-lapata-2008-modeling}。网格以句子为行、共指实体为列，格值为S（主语）、O（宾语）、X（其他成分）或--（不出现）；同一句多次出现同一实体时，按S优先于O、O优先于X取值。每列从上到下的相邻格值构成长度为2的转移，如SS、S--、OS，共16种；统计全文各类转移所占的比例，得到特征向量$\bm\phi(D)\in\mathbb R^{16}$。与中心理论不同，网格不规定哪种过渡更好，而是从数据中学习每种转移的权重。

对三句材料，把“他”并入小王、“这本书”并入书，网格为
\[
  \begin{array}{c|cc}
     & \text{小王} & \text{书}\\ \hline
   A & \mathrm S & \mathrm O\\
   B & \mathrm S & \mathrm O\\
   C & - & \mathrm S
  \end{array}.
\]
C是被动句，书在语义上是受事，句法上是主语，格值取S。两列各有两次转移，共4次，SS、S--、OO、OS各占$1/4$。

学习用成对排序：原文$D$应比它的扰动版本$D'$得分高，设$f(D)=\bm w^{\top}\bm\phi(D)$，排序SVM的目标函数为
\[
  \min_{\bm w}\ \frac{\mu}{2}\lVert\bm w\rVert_2^2+\sum_{(D,D')}\max\bigl(0,1-f(D)+f(D')\bigr),
\]
其中$\mu$为正则化系数。设教学权重$w_{\mathrm{SS}}=2$、$w_{\mathrm{OO}}=1$、其余为0，原顺序得分$0.75$；ACB产生S--、--S、OS、SO四种转移，得分0，分差$0.75$，铰链损失为$0.25$。网格依赖共指结果：若漏掉“他”与小王的共指，网格变成三列，SS消失，OO只占$1/6$，原文得分降为$1/6$，共指错误由此传入连贯性判断。实体网格只捕捉实体延续，一段始终围绕同一实体、却前后矛盾的文字，网格无法区分。Nguyen与Joty把网格的离散特征换成可学习的表示：每个格值S、O、X、--映射为嵌入向量，在每一列（即每个实体的角色序列）上做一维卷积以捕捉更长的转移模式，池化后打分，仍以原文优于扰动的成对排序损失训练\scite{nlp-analysis-nguyen-joty2017}。卷积窗口覆盖多句，弥补了长度为2的转移只看相邻句的局限，但输入仍是实体网格，共指错误同样会传入。

\subsubsection{句对衔接评分与束搜索排序}
\label{sec:nlp-analysis-pair-ordering}

实体网格给整篇打分，要用它恢复句序就得枚举排列。另一种思路直接学习“哪句适合接在哪句后面”：训练有方向的句对评分$s(a,b)$，正样本来自原文中的相邻句，负样本来自不相邻句或颠倒的句对\scite{nlp-analysis-chen2016ordering}。再学习起始分数$s_0(a)$，表示句子$a$适合作为开头。对排列$\pi=(\pi_1,\ldots,\pi_m)$，总分为
\[
  S(\pi)=s_0(\pi_1)+\sum_{t=2}^{m}s(\pi_{t-1},\pi_t).
\]
求最高分排列等价于在句子构成的完全图上找一条最大权哈密顿路径，一般是NP难的，实践中用束搜索近似（第~\ref{sec:nlp-generation-beam}小节），搜索状态为“已选序列及剩余句子集合”。

取$s_0(A)=2$、$s_0(B)=0$、$s_0(C)=-1$，$s(A,B)=3$、$s(A,C)=4$、$s(B,C)=2$、$s(B,A)=-1$、$s(C,A)=-3$、$s(C,B)=-2$（教学构造）。束宽为2时，第一步保留A（2分）与B（0分）；扩展后AC得6、AB得5、BC得2、BA得$-1$，保留AC与AB；第三步各自只能接剩下的句子，ACB得$6-2=4$，ABC得$5+2=7$，输出ABC。束宽为1时，第二步只保留AC，最终得到ACB。穷举全部6个排列可确认7分是最大值。这个例子说明，句对评分只看相邻两句，局部最高分的选择不保证整条排列最优；束宽越大越接近最优，计算量也随之增加。

句对评分还有另一种用法：不判断“是否相邻”，而判断“$a$是否在$b$之前”。Prabhumoye等用BERT对每一对句子分类先后，得到一张有向图，再对它做拓扑排序得到全序\scite{nlp-analysis-prabhumoye2020}。先后关系比相邻关系更容易从内容判断（例如“读完”应早于“归还”，中间可以隔着别的句子），而且拓扑排序是线性时间的精确算法；代价是成对分类可能互相矛盾而形成环，需要先按置信度删边，并且$O(m^2)$个句对都要过一次编码器。

\subsubsection{指针式句子排序}
\label{sec:nlp-analysis-pointer-ordering}

句对评分只依赖前一句，“B接在A后面合不合适”与A之前已经说过什么无关。指针网络让每一步的选择依赖已选的全部历史\scite{nlp-analysis-gong2016pointer}。编码器先把每个句子编码为$\bm h_j$，再对全部句子做整体编码（如以句子表示为输入的LSTM或自注意力层），得到段落表示作为解码器初始状态；解码器把已选句子依次输入，维护状态$\bm z_t$，在第$t$步对每个剩余句子打分：
\[
  g_{t,j}=\bm v^{\top}\tanh(\bm W\bm h_j+\bm U\bm z_t),\qquad
  p(\pi_t=j\mid\pi_{<t})=\operatorname{softmax}_{j\notin\pi_{<t}}(g_{t,j}),
\]
已选句子的分数设为$-\infty$。训练时把参考顺序作为解码输入（教师强制），对每一步取交叉熵。

设第一步对A、B、C的概率为$0.6,0.3,0.1$，选A；第二步屏蔽A，对B、C的概率为$0.75,0.25$，选B；第三步只剩C。ABC的路径概率为$0.6\times0.75=0.45$，训练损失$-\log0.45\approx0.799$；ACB的路径概率为$0.6\times0.25=0.15$。训练时解码器总是看到正确的历史，推断时则看到自己的预测，首句选错会影响此后每一步的分数，这是第~\ref{sec:nlp-analysis-paradigm-transition}小节所说的暴露偏差在句子排序中的表现。

\paragraph{排序评价}
排序结果用两个指标评价。Kendall τ衡量预测排列与参考排列的一致程度\scite{nlp-analysis-lapata2006}。设$m$个句子有$\binom m2$个句子对，其中在预测中相对次序与参考相反的对数为$I$（逆序对数），则
\begin{equation}
  \tau=1-\frac{2I}{\binom m2},
  \label{eq:nlp-analysis-kendall-tau}
\end{equation}
取值在$-1$到1之间，完全一致为1，完全颠倒为$-1$。完全匹配率（perfect match ratio, PMR）是预测排列与参考完全一致的文档所占比例。对ACB，三个句子对中只有(B,C)颠倒，$I=1$，$\tau=1-2/3=1/3$。若测试集有两篇文档，一篇完全恢复，一篇输出ACB，则平均τ为$(1+1/3)/2\approx0.667$，PMR为$1/2$。τ对局部交换给部分分，PMR只认完全正确，两者结合可以区分“大体对、局部错”和“完全错”。两个指标都假设参考顺序唯一，存在多种合理顺序时会低估模型。

表~\ref{tab:nlp-analysis-coherence-methods}汇总本小节的方法。

\begin{table*}[htbp]
  \centering
  \small
  \begin{tabularx}{\linewidth}{@{}l XXX@{}}
    \toprule
    方法 & 解决的问题 & 核心机制 & 新的代价或局限\\
    \midrule
    中心理论 & 描述局部连贯 & 回指中心与首选中心的过渡类型及偏好次序 & 不可学习；不刻画事件顺序\\
    实体网格 & 把实体转移变为可学习特征 & 角色转移比例特征加成对排序损失 & 依赖共指与句法；只看实体延续\\
    句对评分加束搜索 & 直接恢复句序 & 有方向的相邻句评分，束搜索求排列 & 只依赖前一句；搜索近似\\
    指针网络 & 选择依赖全部历史 & 编码全部句子，逐步指针选择剩余句子 & 暴露偏差；长文档成本高\\
    \bottomrule
  \end{tabularx}
  \caption{连贯性与句子排序方法对照。}
  \label{tab:nlp-analysis-coherence-methods}
\end{table*}
